\clearpage\thispagestyle{empty}
\pdfbookmark[0]{Quellen zu den heliozentrischen Vorschriften (1802)}{vorschriften1802-sources}
\begin{center}\textbf{Quellen zu den heliozentrischen Vorschriften (1802)}\end{center}
\noindent Diese KI-Übertragung des deutschen Originals in heutiges Deutsch folgt Carl Friedrich Gauß’ Text in den \textit{Werken}, Band~VI (Göttingen, 1874), S.~87--93. Der Erstdruck erschien im Juniheft 1802 der von Franz Xaver von Zach herausgegebenen \textit{Monatlichen Correspondenz zur Beförderung der Erd- und Himmels-Kunde}, Band~5, S.~540--546, unter dem Namen Dr.~Gauß in Braunschweig. \href{https://gdz.sub.uni-goettingen.de/id/PPN236007467}{Historischer deutscher Text in den Werken}. \href{https://www.digitale-sammlungen.de/en/view/bsb10538597?page=556}{Erstdruck}.
\par\smallskip\noindent Im Erstdruck steht \(V\) für die Länge der Sonne, in den \textit{Werken} \(\odot\). Im Rechenbeispiel \(6^{0}\) steht im Erstdruck hinter dem Komplement von \(\log\cos i\) ein \(n\) (S.~544). Die \textit{Werke} geben \(0.0075025\) ohne dieses Zeichen (S.~91). Die Übertragung folgt der Lesart der \textit{Werke}.
\par\smallskip\noindent Ein Nachdruck steht in Anhang IV, S.~42--45, zu \href{https://gdz.sub.uni-goettingen.de/id/PPN60187157X}{Carl Haases deutscher Übersetzung der \textit{Theoria motus}} (Hannover, 1865). Gauß’ verwandte Darstellung und seinen Rückverweis auf den Zeitschriftenaufsatz findet man in Artikel~74, S.~89--91, in \href{https://www2.karlin.mff.cuni.cz/~tuma/Aplikace21/theoryofmotionof00gausuoft.pdf#page=113}{Charles Henry Davis’ englischer Übersetzung der \textit{Theoria motus}} (Boston, 1857).
\par\smallskip\noindent Übertragung und Satz wurden mit KI erstellt.
\clearpage
